\section{Conclusion}

As already stated in the introduction, the goal of this project was to crate a system that can help users navigate through their environment by providing spatial guidance.
In this work we have shown a 3D spatial scene description system that leverages Gemini VLM to translate visual and depth information into a way to explane the environment to the user.
With the solution crated it is possible to get the location or height of objects, as well as further orientation assistance.

Returning to the original research question, we can conclude that the work in this field is still in its infancy and even thoughout the time of working on this project there were multible relesases that changed our apprach.
None of the models that are used in this project were available at the start of this research project.
This also shows that there is still work to do and that the field is still evolving with could potentialy outdate this solution in the following months.
Still the results showen in the evaluation section reveal that it is possible to get accurate and consistend navigation and spatial guidance information by using the Gemini Pro Model and Gemini Live API.
This way it is not only possible to get accurate information, but the live api also allows for a easy to implement way of creating a real time conversation with the user.

\subsection{Future Work}

As showen in the discussion section, there are still limitations of the current solution hat can be fixed in the future.
The two most important ones are Model and Performance enhancements and System integration:

\begin{itemize}
    \item \textbf{Model and Performance Enhancements:} One aspekt in order to improve the performance would be to futher use example in- and outputs to get better json quality and results. Additionally, add ensembling techniques to the system could also help. This would mean to use bach requests in order to get multiple results for the same input and then use a voting system to decide witch result to pick. Another would be to seperate the segementation and formatting in two different models. The one that does the segementation could use a higher temperature while the one for formatting would use a lower temperature. This way the segmentation model could be more creative in its outputs, while the formatting model would be more strict and consistent in its outputs. Another idea would be to introduce a confidence score for the segementation masks, which would allow the user to see how confident the model is in its segmentation. This could help the user to understand the quality of the segmentation and make better decisions based on that.
    \item \textbf{System Integration:} Another aspect that can be enhanced is the implementation of the system in a real world scenario. This would include the integration of the system into a moble app with would come with its own challenges. For exampele, the system would need to use Websocket connections to the Gemini Live API in order to provice real time conversations.
\end{itemize}
